\chapter{Introduction}

\section{Problem Statement}

As the recent earthquake in Thailand caused significant structural damage, particularly to high-rise buildings which led to widespread concrete cracking and breakage. Therefore, the need for efficient and reliable damage assessment has become more critical than ever. Structural defects such as cracks and spalling not only undermine the integrity of infrastructure but also pose serious risks to public safety if not detected and repaired promptly.

Traditional inspection methods rely heavily on manual labor which are time-consuming and are prone to human error making them impractical for large scale assessments especially in post-disaster scenarios. By that we want to enhance the current method by using modern technology from this project instead.

This project proposes an AI-powered defect detection system that leverages deep learning to automate the identification and localization of structural damage from images.By utilizing pre-trained models such as ResNet, EfficientNet and MobileNetV3. The system classifies images as either cracked or non-cracked and applies object detection techniques for precise localization. This approach not only improves the speed and accuracy of structural assessments but also supports engineers and decision-makers in prioritizing repairs, ultimately enhancing infrastructure resilience after natural disasters.


\section{Project Scope and Goals}

The goal of this project is to enhance structural health monitoring through automated crack detection, reducing manual inspection and improving safety. To achieve this, the project will:

\begin{enumerate}
	\item Develop a classification model (ResNet, EfficientNet and MobileNetV3) in order to distinguish between cracked and non-cracked surfaces with high accuracy.
	\item Implement object detection and segmentation using Faster R-CNN for precise localization of cracks.
	\item Preprocess and enhance data, including image augmentation and feature extraction for improved model performance.
	\item Train and evaluate models using relevant performance metrics such as accuracy, precision, recall, F1-score, AUC-ROC and Confusion Matrix.
	\item Deploy the best-performing model via Flask/FastAPI, Docker and cloud platforms (AWS, GCP, Azure or Hugging Face Spaces) for real-world applications.
	\item Explore optional CI/CD integration, including automated model retraining, performance monitoring, and deployment pipelines.
\end{enumerate}

This project will contribute to efficient and scalable crack detection, improving infrastructure maintenance and safety through AI-driven automation.

\section{Data Description}

\section{Keywords}



\clearpage